\chapter{Spontaneous abortion (miscarriage)}
\index{Spontaneous abortion (miscarriage)}

\section*{Definition and Overview}
Miscarriage, clinically referred to as spontaneous abortion, is the involuntary loss of a pregnancy before the 20th week of gestation (or when the fetus weighs less than 400 grams, as defined in local clinical practice). Miscarriage is the most common complication of early pregnancy, occurring in various clinical presentations including threatened, inevitable, complete, incomplete, missed, or septic miscarriage. Recurrent miscarriage is defined as the loss of three or more consecutive pregnancies. The vast majority of early miscarriages are due to chromosomal abnormalities that occur at conception. Spontaneous pregnancy loss is a profound physical and emotional experience requiring sensitive clinical management, supportive counselling, and patient education.

\section*{Epidemiology}
Miscarriage is highly prevalent, affecting approximately 10\% to 20\% of all clinically recognised pregnancies. It is estimated that up to 50\% of all conceptions end in miscarriage, with many occurring before the woman is aware of the pregnancy (chemical pregnancy). In many countries, around 100,000 miscarriages are recorded annually. The risk of miscarriage increases significantly with advancing maternal age: from approximately 10\% to 15\% in women under 30 to over 50\% in women over 45, primarily due to age-related declines in oocyte quality. Paternal age over 40 is also associated with an increased risk. Approximately 1\% of couples experience recurrent miscarriage.

\section*{Aetiology and Pathophysiology}
Miscarriage is caused by a variety of genetic, anatomical, endocrine, or immunological factors:
\begin{itemize}
    \item \textbf{Chromosomal abnormalities:} Accounts for 50\% to 60\% of first-trimester miscarriages, involving spontaneous errors in cell division (aneuploidies, trisomies, monosomy X, or polyploidies) unrelated to maternal actions.
    \item \textbf{Maternal age:} Advancing age increases the rate of chromosomal non-disjunction during oogenesis.
    \item \textbf{Endocrine factors:} Inadequate progesterone support (luteal phase defect), uncontrolled diabetes mellitus, thyroid disease, or polycystic ovary syndrome (PCOS).
    \item \textbf{Anatomical abnormalities:} Uterine septae, submucosal fibroids, cervical insufficiency, or intrauterine adhesions (Asherman syndrome).
    \item \textbf{Immunological conditions:} Antiphospholipid syndrome (APS), which causes placental thrombosis and vascular dysfunction, leading to recurrent pregnancy loss.
    \item \textbf{Lifestyle and environmental exposures:} Smoking, heavy alcohol consumption, high caffeine intake, drug abuse, or exposure to specific environmental toxins.
\end{itemize}

\section*{Clinical Features}
The symptoms and signs of miscarriage depend on the gestational age and clinical subtype:
\begin{itemize}
    \item \textbf{Vaginal bleeding:} Spotting, light bleeding, or heavy bleeding with clots, which is the most common presenting symptom of early pregnancy loss.
    \item \textbf{Abdominal cramping:} Pelvic pain, cramping, or lower back ache, which may range from mild (similar to a period) to severe.
    \item \textbf{Passage of tissue:} Expulsion of greyish-pink tissue or gestational sac from the vagina.
    \item \textbf{Loss of pregnancy symptoms:} Sudden regression of breast tenderness, morning sickness, or fatigue.
    \item \textbf{Systemic signs (Septic miscarriage):} Fever, chills, lower abdominal tenderness, and foul-smelling vaginal discharge.
\end{itemize}

\section*{Differential Diagnosis}
Early pregnancy bleeding and pain require careful evaluation to distinguish miscarriage from other conditions:
\begin{itemize}
    \item \textbf{Ectopic pregnancy:} Implantation of the embryo outside the uterine cavity (most commonly in the fallopian tube), a life-threatening emergency presenting with bleeding and localised pelvic pain.
    \item \textbf{Molar pregnancy (hydatidiform mole):} A gestational trophoblastic disease presenting with heavy bleeding, hyperemesis, abnormally high hCG levels, and a "snowstorm" pattern on ultrasound.
    \item \textbf{Implantation bleeding:} Light spotting occurring around the time of the missed period, which is benign and common.
    \item \textbf{Cervical or vaginal pathology:} Bleeding from cervical polyps, ectropion, or local infections, distinguished by speculum examination.
    \item \textbf{Urinary tract infection (UTI):} Presents with lower abdominal pain and dysuria, but no vaginal bleeding.
\end{itemize}

\section*{When to Seek Medical Care}
Any pregnant individual who experiences vaginal bleeding, spotting, or pelvic pain in the first 20 weeks should contact their obstetrician or visit an early pregnancy clinic.

\textbf{Contact your local emergency services for:}
\begin{itemize}
    \item \textbf{Severe, sharp, or one-sided pelvic pain:} Especially if accompanied by shoulder tip pain, dizziness, or fainting (indicates ruptured ectopic pregnancy).
    \item \textbf{Haemorrhage:} Heavy bleeding that saturates more than two sanitary pads per hour for two consecutive hours.
    \item \textbf{Signs of septic miscarriage:} High fever, chills, rapid heart rate, confusion, or severe abdominal tenderness.
    \item \textbf{Loss of consciousness or shock:} Severe lightheadedness or collapse.
\end{itemize}
Other reasons to see a doctor include passage of tissue or clots, mild pelvic cramping that worsens, or persistent brown discharge.

\section*{Diagnosis and Investigations}
Establishing the status of the pregnancy involves a combination of clinical, biochemical, and imaging assessments:
\begin{itemize}
    \item \textbf{Transvaginal ultrasound (TVS):} The primary diagnostic tool. A missed or incomplete miscarriage is confirmed by the absence of a fetal heartbeat in an embryo with a crown-rump length (CRL) of 7 mm or more, or a mean sac diameter (MSD) of 25 mm or more with no embryo.
    \item \textbf{Serum beta-hCG measurements:} Sequential blood tests taken 48 hours apart to assess whether hCG levels are rising normally (doubling every 48 hours in early viable pregnancy) or falling.
    \item \textbf{Maternal blood group and Rhesus status:} Crucial to determine the need for anti-D immunoglobulin injection in Rhesus-negative women to prevent isoimmunisation.
    \item \textbf{Full blood count (FBC):} Obtained to assess the severity of blood loss and screen for infection.
    \item \textbf{Histopathological examination:} Analysis of expelled tissue to confirm intrauterine pregnancy and rule out gestational trophoblastic disease.
\end{itemize}

\section*{Management}
Management options for miscarriage include expectant, medical, or surgical care:
\begin{itemize}
    \item \textbf{Expectant management:} Waiting for the pregnancy tissue to pass naturally, suitable for stable patients in the first trimester without signs of infection or excessive bleeding.
    \item \textbf{Medical management:} Administration of misoprostol (a prostaglandin analogue, sometimes preceded by mifepristone) to stimulate uterine contractions and expel the remaining tissue, offering a predictable schedule without surgery.
    \item \textbf{Surgical management:} Dilatation and curettage (D\&C) or suction curettage, performed under local or general anaesthesia. Indicated for heavy bleeding, infection (septic miscarriage), patient preference, or failed medical management.
    \item \textbf{Anti-D immunoglobulin:} Administered to Rhesus-negative women within 72 hours of a miscarriage to prevent maternal sensitisation in future pregnancies.
    \item \textbf{Supportive care:} Access to bereavement support, professional psychological counselling, and organisations such as Sands or the Pink Elephants Support Network.
\end{itemize}

\section*{Complications}
Potential complications associated with miscarriage include:
\begin{itemize}
    \item \textbf{Severe haemorrhage:} Heavy bleeding that may require blood transfusion or emergency surgical evacuation.
    \item \textbf{Septic miscarriage:} A severe uterine infection that can progress to pelvic inflammatory disease, peritonitis, or life-threatening septic shock.
    \item \textbf{Retained products of conception (RPOC):} Placental or fetal tissue left in the uterus, causing persistent bleeding, pain, and infection.
    \item \textbf{Asherman syndrome:} Intrauterine adhesions or scarring that can develop after curettage, potentially leading to future infertility or menstrual issues.
    \item \textbf{Psychological morbidity:} High incidence of clinical depression, anxiety, post-traumatic stress disorder (PTSD), and prolonged grief.
\end{itemize}

\section*{Prognosis and Outcomes}
The physical prognosis following a single miscarriage is excellent, with the vast majority of women recovering fully without any long-term physical sequelae. Over 80\% to 85\% of women who experience a single miscarriage go on to have a successful, healthy pregnancy in the future. The prognosis remains favourable even for couples experiencing recurrent miscarriage, with up to 60\% to 70\% eventually conceiving successfully, especially when underlying causes (such as APS or uterine anomalies) are diagnosed and treated. Psychological recovery varies, and adequate emotional support is vital for long-term well-being.

\section*{Prevention and Self-Care}
While early miscarriage cannot usually be prevented, patients can implement self-care to support recovery and future pregnancies:
\begin{itemize}
    \item \textbf{Manage pain and bleeding at home:} Use sanitary pads instead of tampons during the bleeding phase to reduce the risk of uterine infection.
    \item \textbf{Observe pelvic rest:} Avoid sexual intercourse and douching until all bleeding has stopped and as advised by your doctor.
    \item \textbf{Optimise lifestyle factors:} Quit smoking, avoid alcohol, maintain a healthy weight, and limit caffeine intake when planning a future pregnancy.
    \item \textbf{Take prenatal vitamins:} Begin taking folic acid (at least 400 micrograms daily) before trying to conceive again to reduce neural tube defect risks.
    \item \textbf{Allow time for emotional healing:} Do not rush into trying to conceive immediately; wait until both physical recovery is complete and emotional readiness is reached.
\end{itemize}

\section*{Sources and Further Reading}
The following reputable organisations provide further information on miscarriage:
\begin{itemize}
    \item Pink Elephants Support Network. \textit{Support, resources, and validation for early pregnancy loss.} https://www.pinkelephants.org.au/
    \item Red Nose Australia (incorporating Sands). \textit{Miscarriage, stillbirth, and newborn death support.} https://www.rednose.org.au/
    \item Better Health Channel. \textit{Spontaneous abortion (miscarriage): symptoms and recovery.} https://www.betterhealth.vic.gov.au/
    \item Mayo Clinic. \textit{Miscarriage: symptoms, causes, and diagnosis.} https://www.mayoclinic.org/
    \item NHS. \textit{Miscarriage: overview, treatment, and support services.} https://www.nhs.uk/
\end{itemize}
